\documentclass[12pt]{ximera}
\title{Homework 6: Bayes' Theorem}
\begin{document}
\begin{abstract}
Section 7.6: translating a word problem into conditional-probability notation,
choosing the tree diagram that matches the information given, and applying Bayes'
Theorem --- including the classic medical-testing example where a $90\%$ sensitive
test still gets it wrong a third of the time.
\end{abstract}
\maketitle

\section*{Section 7.6: Bayes' Theorem}

\begin{problem}
Suppose after an election, it is determined that only $60\%$ of eligible voters
actually voted ($V$). Among those who voted, $35\%$ were registered Republican
($R$), $30\%$ Democrat ($D$), $25\%$ Independent ($I$), and $10\%$ other ($O$);
and among those who didn't vote ($\olsi{V}$), $20\%$ were registered Republican,
$25\%$ Democrat, $30\%$ Independent, and $25\%$ other.

\begin{prompt}
\textbf{(a)} Using the set abbreviations given above, translate every probability
from this problem into mathematical notation.

For example, ``only $60\%$ of eligible voters actually voted'' translates to
$P(V)=60\%$, or $0.6$.

\begin{freeResponse}
\end{freeResponse}

\begin{feedback}
The overall voting rate, and its complement:
\[ P(V)=0.60, \qquad P\left(\olsi{V}\right)=0.40. \]

The party breakdowns are all \emph{conditional}, because each was reported within a
group. Among those who voted:
\[
P(R\mid V)=0.35,\quad P(D\mid V)=0.30,\quad P(I\mid V)=0.25,\quad P(O\mid V)=0.10.
\]
Among those who did not vote:
\[
P\left(R\mid\olsi{V}\right)=0.20,\quad
P\left(D\mid\olsi{V}\right)=0.25,\quad
P\left(I\mid\olsi{V}\right)=0.30,\quad
P\left(O\mid\olsi{V}\right)=0.25.
\]
Notice that each group of four conditional probabilities adds to $1$.
\end{feedback}
\end{prompt}

\begin{prompt}
\textbf{(b)} Which tree diagram is appropriate for the information given above?

\begin{multipleChoice}
\choice{Tree \#1: the first split is by party, and the second split is whether they voted}
\choice[correct]{Tree \#2: the first split is whether they voted, and the second split is by party}
\end{multipleChoice}

\begin{hint}
A tree's second-generation branches carry conditional probabilities of the form
$P(\text{second}\mid\text{first})$. Which way around were the percentages in the
problem reported?
\end{hint}

\begin{feedback}
Tree \#2. The problem gives $P(V)$ first, and then reports party membership
\emph{within} each voting group --- $P(R\mid V)$, $P(R\mid\olsi{V})$, and so on.
That is exactly the structure of a tree that splits on voting status first and party
second.

Tree \#1 would require the probabilities $P(V\mid R)$, $P(V\mid D)$, and so forth on
its second-generation branches, and those are not given --- in fact finding one of
them is the point of part (c).
\end{feedback}
\end{prompt}

\begin{prompt}
\textbf{(c)} Using Bayes' Theorem, determine the probability that an eligible voter
who was registered Independent actually voted in the election. That is, determine
$P(V\mid I)$.

Give your answer as a fraction in lowest terms: $P(V\mid I)=\answer{5/9}$

\begin{hint}
Bayes' Theorem here reads
\[
P(V\mid I)=\frac{P(I\mid V)P(V)}{P(I\mid V)P(V)+P\left(I\mid\olsi{V}\right)P\left(\olsi{V}\right)}.
\]
The denominator is just the total probability of being Independent, found by adding
up the two paths through the tree that end in $I$.
\end{hint}

\begin{feedback}
There are two paths through the tree that end at ``Independent'':
\[
P(I\cap V)=P(I\mid V)P(V)=(0.25)(0.60)=0.15,
\]
\[
P\left(I\cap\olsi{V}\right)=P\left(I\mid\olsi{V}\right)P\left(\olsi{V}\right)=(0.30)(0.40)=0.12.
\]
So $P(I)=0.15+0.12=0.27$, and
\[
P(V\mid I)=\frac{P(I\cap V)}{P(I)}=\frac{0.15}{0.27}=\frac{15}{27}=\frac59\approx 0.556.
\]
So about $55.6\%$ of registered Independents voted --- lower than the $60\%$ overall
turnout, which makes sense because Independents were over-represented among
non-voters.
\end{feedback}
\end{prompt}
\end{problem}

\begin{problem}
Consider the following definitions, then use Bayes' Theorem to solve.
\begin{itemize}
    \item The \emph{sensitivity} of a medical test is the probability that the test
      will (accurately) return a positive result, provided whatever it's testing for
      is indeed present.
    \item The \emph{specificity} of a medical test is the probability that the test
      will (accurately) return a negative result, provided whatever it's testing for
      is NOT present.
\end{itemize}

Suppose a blood test has $90\%$ sensitivity and $99\%$ specificity, and suppose it is
known that $2\%$ of people have the particular blood disorder the test is designed to
detect.

\begin{prompt}
\textbf{(a)} What is the probability that someone does indeed have this blood
disorder, given they receive a positive test result?

To the nearest percent: $\answer{65}\%$

\begin{hint}
In symbols, sensitivity is $P(+\mid B)=0.90$ and specificity is
$P\left(-\mid\olsi{B}\right)=0.99$, so the false-positive rate is
$P\left(+\mid\olsi{B}\right)=0.01$. You want $P(B\mid +)$, which reverses the
conditioning --- so Bayes' Theorem applies.
\end{hint}

\begin{feedback}
Write $B$ for having the disorder. We know
\[
P(B)=0.02,\quad P(+\mid B)=0.90,\quad P\left(+\mid\olsi{B}\right)=1-0.99=0.01.
\]
The two paths that end in a positive result are
\[
P(+\cap B)=(0.90)(0.02)=0.018, \qquad
P\left(+\cap\olsi{B}\right)=(0.01)(0.98)=0.0098,
\]
so $P(+)=0.018+0.0098=0.0278$ and
\[
P(B\mid +)=\frac{0.018}{0.0278}=\frac{90}{139}\approx 0.6475,
\]
about $65\%$.

This is the point of the problem: even with a very accurate test, a positive result
means only a $65\%$ chance of actually having the disorder. The disorder is rare, so
the $1\%$ of the healthy $98\%$ who test positive falsely --- nearly $1$ person in
$100$ --- is comparable in size to the $1.8$ people in $100$ who are correctly
flagged.
\end{feedback}
\end{prompt}

\begin{prompt}
\textbf{(b)} What is the probability that someone does NOT have this blood disorder,
given they receive a negative test result?

Give your answer as a percentage rounded to two decimal places:
$\answer[tolerance=0.005]{99.79}\%$

\begin{feedback}
The two paths that end in a negative result are
\[
P\left(-\cap\olsi{B}\right)=(0.99)(0.98)=0.9702, \qquad
P(-\cap B)=(0.10)(0.02)=0.002,
\]
so $P(-)=0.9702+0.002=0.9722$ and
\[
P\left(\olsi{B}\mid -\right)=\frac{0.9702}{0.9722}=\frac{4851}{4861}\approx 0.997943,
\]
that is, about $99.79\%$ --- which rounds to $100\%$ at the nearest whole percent.

Compare this with part (a): the same test is extremely trustworthy when it comes
back negative and only moderately trustworthy when it comes back positive. That
asymmetry comes entirely from the disorder being rare.
\end{feedback}
\end{prompt}
\end{problem}

\end{document}
